% Generated by tools/startup_publication_sync.py.
All current ablations use the pinned vessel-response records and the
final 900 s. Each row below reports the arithmetic mean of eight per-record
RMS scores; ``worst'' is the largest record RMS. These means differ from
root-mean-square pooling used by the engine study.

\subsection{Continuous hard-iron correction}
\begin{table}[t]\centering
\caption{Default-draw hard-iron ablation with current deployed startup.}
\label{tab:startup-hardiron}\footnotesize
\begin{tabular}{@{}lrrrr@{}}\toprule
& \multicolumn{2}{c}{Mean yaw RMS [deg]} & \multicolumn{2}{c}{Worst yaw RMS [deg]} \\
Family & off & on & off & on \\
\midrule
OU--II & 2.137 & 0.581 & 2.198 & 0.727 \\
OU--III & 1.945 & 0.478 & 2.064 & 0.815 \\
TFG & 1.716 & 0.531 & 1.801 & 0.933 \\
\bottomrule\end{tabular}\end{table}

\subsection{TFG feature attribution}
TFG retains both startup policies, allowing matched feature-off replays.
The OU wrappers retain only their deployed proxy startup, so no numerical
comparison to a removed OU startup implementation is reported.
\begin{table}[t]\centering
\caption{TFG default-draw feature ablations. Bias is mean accelerometer-bias RMS.}
\label{tab:tfg-ablation}\footnotesize
\resizebox{\columnwidth}{!}{%
\begin{tabular}{@{}lrrrrr@{}}\toprule
Arm & Yaw [deg] & Roll [deg] & Pitch [deg] & 3-D [m] & Bias [m/s$^2$] \\
\midrule
Deployed & 0.5310 & 0.2923 & 0.0606 & 0.2961 & 0.0488 \\
Hard iron off & 1.7163 & 0.2817 & 0.0609 & 0.2968 & 0.0470 \\
Mag refinement off & 0.7007 & 0.4598 & 0.2211 & 0.2966 & 0.0880 \\
Staged startup & 2.0043 & 0.2683 & 0.0985 & 0.3000 & 0.0482 \\
\bottomrule\end{tabular}}\end{table}

\subsection{Calibration-draw sensitivity}
The five draws are default, 11, 23, 202 and 3033. Each changes the initial
magnetic offset and distortion while holding the paired realization fixed.
For OU--II, mean yaw RMS changes from 3.025 to 1.769 degrees; 3 of 40 pairs worsen.
For OU--III, mean yaw RMS changes from 2.933 to 1.735 degrees; 4 of 40 pairs worsen.
These finite draws do not prove uniform calibration improvement.
